% from ega4-2-fr.tex
\oldpage[IV-2]{55}

\begin{equation}
\tag{4.2.1.3}
\label{IV.2.4.2.1.3-fr}
\xymatrix@C=3.5em@R=2.8em{
K \ar[r] & K\otimes_k L \ar[r] & K\otimes_{K'}L(\mathbf t)\\
K' \ar[u] \ar[r] & K'\otimes_k L \ar[u] \ar[r] & L(\mathbf t) \ar[u]\\
k \ar[u] \ar[r] & L \ar[u] &
}
\end{equation}

Comme $K$ est fidèlement plat sur $K'$, $K\otimes_k L
=K\otimes_{K'}(K'\otimes_k L)$ est fidèlement plat sur $K'\otimes_k L$,
donc $K'\otimes_k L$ est noethérien $(0_{\mathrm I},6.5.2)$~; en outre
$K'\otimes_k L$ s'identifie à un sous-anneau de $K\otimes_k L$~; la trace
sur $K'\otimes_k L$ d'un idéal premier $\mathfrak p$ associé à $K\otimes_k L$
est l'idéal $0$, un élément $\neq0$ de $K'\otimes_k L$ n'étant pas diviseur de
zéro dans $K\otimes_k L$ $(0_{\mathrm I},6.3.4$ et Bourbaki,
\emph{Alg. comm.}, chap.~IV, \S~1, n\textsuperscript{o}~1, cor.~3 de la
prop.~2). Comme en outre $K$ est algébrique sur $K'$, $K\otimes_k L$ est une
algèbre entière sur $K'\otimes_k L$, et les idéaux premiers de $K\otimes_k L$
induisant $0$ sur $K'\otimes_k L$ sont nécessairement sans relation
d'inclusion mutuelle (Bourbaki, \emph{Alg. comm.}, chap.~V, \S~2,
n\textsuperscript{o}~1, cor.~1 de la prop.~1)~; cela prouve la première
assertion (Bourbaki, \emph{Alg. comm.}, chap.~IV, \S~1,
n\textsuperscript{o}~3, cor.~1 de la prop.~7). De plus, le corps résiduel $E$
de $\mathfrak p$ est algébrique sur le corps résiduel de l'idéal $(0)$ de
$K'\otimes_k L$, c'est-à-dire $L(\mathbf t)$~; donc
$\operatorname{deg.tr.}_L E=\operatorname{deg.tr.}_L L(\mathbf t)
=\operatorname{Card}(I)=\operatorname{deg.tr.}_k K$, autrement dit on a la
première relation (4.2.1.1)~; échangeant les rôles de $K$ et $L$, on a la
seconde relation (4.2.1.1), d'où (4.2.1.2).
\end{proof}

\begin{corollary}[4.2.2]
\label{IV.2.4.2.2-fr}
Sous les hypothèses de (3.3.6), si les préschémas $T_{x,y}$ sont localement
noethériens, ils n'ont pas de cycle premier associé immergé.
\end{corollary}

\begin{corollary}[4.2.3]
\label{IV.2.4.2.3-fr}
Sous les hypothèses de (3.3.6) (resp. (3.3.7)), si les $T_{x,y}$ sont
localement noethériens et si $\mathcal F$ et les $\mathcal G_s$ (pour $s\in S$)
(resp. $\mathcal F$ et $\mathcal G$) sont sans cycle premier associé immergé,
il en est de même de $\mathcal F\otimes_S\mathcal G$.
\end{corollary}

Cela résulte de (4.2.2), (3.3.2) et de la démonstration de (3.3.6). En
particulier, comme tout préschéma sur un corps $k$ est \emph{plat sur $k$}, on
a ainsi démontré l'assertion (i) de la

\begin{proposition}[4.2.4]
\label{IV.2.4.2.4-fr}
Soient $k$ un corps, $X$ et $Y$ deux $k$-préschémas localement noethériens
tels que $X\times_k Y$ soit localement noethérien. Supposons de plus $X$ et
$Y$ intègres. Alors~:
\begin{enumerate}
\item[(i)] $X\times_k Y$ est sans cycle premier associé immergé~; chacune des
composantes irréductibles de $X\times_k Y$ domine $X$ et $Y$, et l'ensemble de
ces composantes est en correspondance biunivoque avec l'ensemble des
composantes irréductibles de
$\operatorname{Spec}(R(X)\otimes_k R(Y))$ (autrement dit l'ensemble des
idéaux premiers minimaux de $R(X)\otimes_k R(Y)$), où $R(X)$ et $R(Y)$ sont
les corps des fonctions rationnelles de $X$ et $Y$ respectivement.
\item[(ii)] Si un point maximal $z$ de $X\times_k Y$ est identifié à un idéal
premier minimal $\mathfrak p$ de $R(X)\otimes_k R(Y)$, l'anneau local
$\mathcal O_{X\times_k Y,z}$ est isomorphe à l'anneau de fractions
$(R(X)\otimes_k R(Y))_{\mathfrak p}$. En particulier, si $R(X)$ ou $R(Y)$ est
séparable sur $k$, $X\times_k Y$ est réduit.

